\documentclass[11pt,a4paper]{moderncv}

% moderncv themes
\moderncvcolor{burgundy}
\moderncvstyle[left]{banking} % style options are 'casual' (default), 'classic', 'oldstyle' and 'banking'
\sethintscolumnlength{2.8cm}

% Typical packages
%\usepackage[utf8]{inputenc} % replace by the encoding you are using
%\usepackage[T1]{fontenc} % Output font encoding for international characters
\usepackage{amsmath}
\usepackage{csquotes}
% Language related packages
\usepackage{fontspec}
\usepackage[english]{babel}
\babelprovide[import]{arabic}
\babelfont[arabic]{rm}{vazirmatn-Regular.ttf}
\usepackage{xeCJK}
\setCJKmainfont{NotoSansTC-Regular.otf}
\setCJKsansfont{NotoSansTC-Regular.otf}
\setCJKmonofont{NotoSansTC-Regular.otf}
\usepackage[left=1.5cm,right=1.8cm,top=2.0cm,bottom=2.0cm]{geometry}
\recomputelengths

\usepackage[perpage, para, symbol]{footmisc}
\usepackage[
    backend=biber,
    style=nature,
    maxnames=15,
    url=false,
    eprint=false,
    isbn=false]{biblatex}
\addbibresource{Publication.bib}

% personal data
\name{Amir}{Fathi}
\title{個人履歷}
\mobile{+886 96 3534 132}
\email{fathi0amir@gmail.com}
\homepage{fathi0amir.github.io}
% --------------------------------------
% Begin here
% --------------------------------------
\begin{document}
\makecvtitle
\vspace{-1.0cm}
% ---------------------------------------
% Occupation
% ---------------------------------------
\section{工作經歷}
\cventry{2024--至今}{博士後研究員}
        {國立陽明交通大學應用化學系}
        {台灣新竹}{}{}
\cventry{2023--2024}{博士後研究員}
        {國立清華大學光電研究中心}
        {台灣新竹}{}{}
\cventry{2020--2023}{博士後研究員}
        {中央研究院分子生物研究所}
        {台灣台北}{}{}
%------------------------------------------
% Education
%------------------------------------------
\section{學歷}
%----------------------------
% NCTU
%----------------------------
\cventry{2013--2020}{博士}
        {國立交通大學應用化學系}{台灣新竹}{}{}
%------------------------------------------
%  Semnan University
%------------------------------------------
\cventry{2008--2010}{碩士}
        {塞姆南大學物理系，伊朗塞姆南}{}{}{}
%------------------------------------------
%  Shahid Beheshti
%------------------------------------------
\cventry{2003--2008}
        {學士}{沙希德·貝赫什提大學物理系，伊朗德黑蘭}{}{}{}
% -----------------------------------
% Projects & Research
% -----------------------------------
\section{研究計畫}
\cvitem{2024--至今}
        {從追蹤模擬生物系統之粒子軌跡中生成影像對比度}
\cvitem{2024--至今}
        {建立干涉式單粒子追蹤裝置，
        用於高速粒子追蹤}
\cvitem{2023--2024}
        {結合 off-axis 顯微鏡與瞬態光柵 PL（時間解析度低於 80 fs），
        以研究 TMD 材料}
\cvitem{2020--2023}
        {利用自製雙光子成像技術，
        在虛擬實境系統中
        研究成年斑馬魚的空間神經編碼}
\cvitem{2017--2020}
        {瞬態吸收影像與 SEM 影像共定位，
        以研究光電元件中的電荷轉移異質性}
\cvitem{2014--2017}
        {建構 pump-probe 顯微鏡，
        用於次繞射極限解析度的奈米粒子追蹤}
\cvitem{2017--2018}
        {建構搭載 TE-Cooled PD 之 SS-PL 系統，
        並以 lock-in amplifier 解調，
        使 NIR 區域靈敏度達到可見光區光子計數 PMT 的水準}
\cvitem{2016--2017}
        {設計、模擬、製作、組裝並測試調諧放大電路，
        作為 lock-in amplifier 的低成本替代方案}
\cvitem{2013--2014}
        {鈣鈦礦太陽能電池的飛秒弛豫研究}
\cvitem{2010--2012}
        {以 Rietveld refinement XRD 分析，
        測定以 CVT 法生長之
        ZnS\textsubscript{x}Se\textsubscript{1-x} 單晶的成分值}
\cvitem{2008--2010}
        {以 Chernov 體擴散模型模擬、生長並表徵
        以 CVT 法製備之
        II--VI 族單晶}
%------------------------------------------
%  Experiences
%------------------------------------------
\section{研究經驗}
\cvitem{2024--至今}
        {模擬生物系統中正常與異常擴散，
        適用於單粒子追蹤實驗}
\cvitem{2024--至今}
        {編程高速相機系統，
        實現非阻塞式影像擷取}
\cvitem{2023--2024}
        {編寫 PIMAX3 驅動程式；建構 off-axis 顯微鏡；
        以 FROG 進行超快雷射脈衝表徵}
\cvitem{2020--2023}
        {建構雙光子雷射掃描顯微鏡，
        整合虛擬實境環境，
        並透過扭矩／力感測器偵測虛擬游泳事件}
\cvitem{2018--2020}
        {操作掃描電子顯微鏡}
\cvitem{2013--2020}
        {架設超快 pump-probe 雷射掃描顯微鏡}
\cvitem{2013--2020}
        {光電元件的超快雷射光譜與顯微研究}
\cvitem{2008--2011}
        {在塞姆南大學進行CVT晶體生長最佳化}
%------------------------------------------
%  Rewards and Honors
%------------------------------------------
\section{獲獎與榮譽}
\cvitem{2021}
        {獲頒「2021 Academia Sinica Postdoctoral Research Scholars」，
        獲得兩年期研究經費補助}
\cvitem{2020}
        {優秀論文獎：``A Direct Mapping
        Approach to Understand Carrier Relaxation
        Dynamics in Varied Regions of a
        Polycrystalline Perovskite Film''}
\cvitem{2020}
        {優秀論文獎：``Label-Free Optical
        Microscope Based on a Phase-Modulated Femtosecond
        Pump--Probe Approach with Sub-diffraction Resolution''}
\cvitem{2018}
        {優秀論文獎：``Slow surface
        passivation and crystal relaxation with additives to
        improve device performance and durability for tin-based
        perovskite solar cells''}
\cvitem{2013--2020}
        {獲得國立交通大學博士班獎學金及學費減免}
\cvitem{2008--2010}
        {獲得政府學術獎學金（碩士班）}
\cvitem{2003--2007}
        {獲得政府學術獎學金（學士班）}
%------------------------------------------
% Conference Attended
%------------------------------------------
\section{參加學術會議}
\cvitem{2018年12月}{\textbf{台灣光電學會年度大會}，
        \newline NCTU（台南校區），台灣}
\cvitem{2012年3月}{\textbf{第四屆奈米結構國際研討會（ICNS4）}，\newline 伊朗基什島}
\cvitem{2010年9月}{\textbf{伊朗年度物理研討會}（伊朗物理學會主辦），
        \newline 布阿里辛那大學，伊朗哈馬丹}
\cvitem{2009年1月}{\textbf{量子計算與量子資訊處理及量子計算實驗面向研討會}，
        \newline 沙希德·貝赫什提大學（伊朗）與近畿大學（日本）聯合主辦，伊朗德黑蘭}
%------------------------------------------
%  Technical Skills
%------------------------------------------
\section{技術能力}
    \subsection{電腦技能}
    \cvlistitem{\small \raggedright MATLAB、LabVIEW、R、Python、C\# 及 C++ 程式設計}
    \cvlistdoubleitem{\small \raggedright 使用 Blender 進行網格（3D）建模}
                    {\small \raggedright 使用 UE5 及 Panda3D 進行 VR 開發}
    \cvlistdoubleitem{\small \raggedright 使用 SolidWorks 進行 CAD 設計及 3D 列印}
                    {\small \raggedright 使用 \TeX、Typst 及 Pandoc 排版}
    \cvlistdoubleitem{\small \raggedright 使用 ImageJ、MATLAB、Napari 進行影像處理}
                    {\small \raggedright 使用 FullProf Suite 進行 XRD 資料分析}
    \subsection{科學儀器}
    \cvlistitem{\small \raggedright 以 PD、PMT、APD、EMCCD 及 iCCD 為偵測器的科學 DAQ 程式設計}
    \cvlistdoubleitem{\small \raggedright 多區箱式與管式爐的校準與維護}
                    {\small \raggedright 穩態光致發光及紫外-可見光測量}
    \cvlistdoubleitem{\small \raggedright 建構 (fs and ns) 瞬態吸收光譜系統}
                    {\small \raggedright 建構 fs pump-probe 顯微鏡及其影像處理流程}
    \cvlistdoubleitem{\small \raggedright 被動/主動電子濾波器的設計、模擬與製作}
                     {\small \raggedright 具有操作掃描電子顯微鏡（SEM）的認證}
%------------------------------------------
%  Languages
%------------------------------------------
\section{語言能力}
\cvitemwithcomment{波斯語}{母語}{\hspace{-0.8em}\dotfill \foreignlanguage{arabic}{پارسی}}
\cvitemwithcomment{英語}{流利 {\footnotesize（iBT 成績：95）}}{\hspace{-0.8em}\dotfill English}
\cvitemwithcomment{西班牙語}{中級}{\hspace{-0.8em}\dotfill Español}
\cvitemwithcomment{德語}{中級}{\hspace{-0.8em}\dotfill Deutsch}
% \cvitemwithcomment{摩斯密碼}{中級}{\hspace{-0.8em}\dotfill \Large{\textbf{ -\hspace{0.1mm}- -\hspace{0.1mm}-\hspace{0.1mm}- .-. ... . -.-. -\hspace{0.1mm}-\hspace{0.1mm}- -.. .}}}
\cvitemwithcomment{中文}{中級}{}

% \newpage
% \nocite{*}
% \printbibliography[title={學術著作}]
\vfill
\raggedleft
\footnotesize{如需推薦信，歡迎索取\\
最後更新：\textit{\today}}
\end{document}
